\documentclass[11pt, headheight=30pt]{scrartcl}
\usepackage[white, nologo, hw]{darkWhite}
\usepackage{mathrsfs}

\newcommand{\BBB}{\mathscr{B}}
\DeclareMathOperator{\cov}{cov}
\DeclareMathOperator{\var}{var}

\title{18.675 Problem Set 4}
\begin{document}
% PROBLEMS SET 4 FOR 18.675 (FALL 2023)
% Problem 1. Let X = (X1, . . . , Xn) be a random variable, with all components in L2(P). The
% covariance matrix var (X) = (ci,j : 1 ≤ i, j ≤ n) of X is defined by ci,j = cov (Xi, Xj ). Show that
% var(X) is a non-negative definite matrix.
\section{}
\begin{prob}
Let $X = (X_1, \dots, X_n)$ be a random variable,
with all components in $L^2(P)$. The covariance
matrix $\var(X) = (c_{i,j} : 1 \leq i, j \leq n)$
of $X$ is defined by $c_{i,j} = \cov(X_i, X_j)$.
Show that $\var(X)$ is a non-negative definite matrix.
\end{prob}

Well, for any particular vector $v = (v_1, \dots, v_n)$,
we have
\begin{align*}
    v^T \var(X) v &= \sum_{i,j} v_i v_j c_{i,j} \\
    &= \sum_{i,j} v_i v_j \cov(X_i, X_j) \\
    &= \cov\left( \sum_i v_i X_i, \sum_j v_j X_j \right) \\
    &= \var\left( \sum_i v_i X_i \right) \\
    &\geq 0
\end{align*}
The third line is true because covariance is linear
in each random variable; this is clear by starting at the definition.
The fourth line is true by definition of variance. The
fifth line is true because squares of real numbers
are nonnegative, and the integral of a non-negative
function is non-negative.


% Problem 2. For a finite Borel measure μ on R show that, if ∫
% R |x|kdμ(x) < ∞ for k ∈ N,
% then the Fourier transform ̂ μ of μ has a kth continuous derivative, which at 0 is given by
% ̂
% μ(k)(0) = ik
% ∫
% R
% xkdμ(x).
\section{}
\begin{prob}
For a finite Borel measure $\mu$ on $\RR$ show that,
if $\int_\RR \abs{x}^k d\mu(x) < \infty$ for $k \in \NN$,
then the Fourier transform $\hat{\mu}$ of $\mu$ has a
$k$th continuous derivative, which at $0$ is given by
\[
    \hat{\mu}^{(k)}(0) = i^k \int_\RR x^k d\mu(x)
\]
\end{prob}

First off, note that on finite measure spaces, $L^p \subset L^q$
whenever $p \leq q$ by Holder's inequality;
\begin{align*}
    \int_\RR \abs{x}^p d\mu(x) &\leq \left(\int_\RR \abs{x}^q d\mu(x)\right)^{p/q} \left(\int_\RR 1 d\mu(x)\right)^{1 - p/q} \\
    &= \left(\int_\RR \abs{x}^q d\mu(x)\right)^{p/q} \mu(\RR)^{1 - p/q} \\
    &< \infty
\end{align*}
So we also know that $\int_\RR \abs{x}^m d\mu(x) < \infty$ for all $m \leq k$.

% The Fourier transform of $\mu$ is defined as
% \[
%     \hat{\mu}(t) = \int_{\RR} e^{itx} d\mu(x)
% \]
% for $t \in \RR$.

Now, let's use induction to show that
\[
    \hat{\mu}^{(m)}(t) = i^m \int_E x^m e^{itx} d\mu(x)\tag{$\star$}
\]
for all $m \leq k$; this clearly suffices (we know $\hat{\mu}^{(m)}$
is continuous because Fourier transforms are continuous, and we can
plug in $0$ for the answer).
The base case is
\[
    \hat{\mu}(t) = \int_{\RR} e^{itx} d\mu(x)    
\]
which is true by definition.

Let $f_m(t,x) = i^m x^me^{itx}$ be the desired integrand. Let $E = \RR$ for clarity.
Observe
\[
    \frac{\pa}{\pa t} f_m(t,x) = f_{m+1}(t,x).
\]

For all $m\geq 0$, we satisfy:
\begin{itemize}
    \item For all $t$, we can integrate $f_m(t,\bullet)$ over $E$. Indeed,
    \[
        \int_E \abs{f_m(t,x)} d\mu(x) \leq \int_E \abs{x}^k d\mu(x) < \infty.
    \]
    \item For all $x$, we can differentiate $f_m(\bullet,x)$ with respect
    to $t$ by construction.
    \item For all $x\in E$ and all $t\in U$,
    \[
        \abs{ \frac{\partial}{\partial t} f_m(t,x) } \leq \abs{f_{k+1}(t,x)}    
    \]
    (in fact they are equal) and the RHS is integrable over $E$.
\end{itemize}
Thus, theorem $3.5.1$ applies, and we can swap the derivative with
the integral; for all $m$,
\begin{align*}
    \frac{\pa}{\pa t} \hat{\mu}^{(m)}(t) &= \frac{\pa}{\pa t} \int_E f_m(t,x) d\mu(x) \\
    &= \int_E \frac{\pa}{\pa t} f_m(t,x) d\mu(x) \\
    &= \int_E i^{m+1} f_{m+1}(t,x) d\mu(x)
\end{align*}
as desired.

% Plugging in $t = 0$, we achieve the desired result.

% I guess I should show that $\star$ is continuous. This is kind of obvious;
% give me some $\eps > 0$. Using the stupid fact that $\abs{e^{i\theta} - 1} < \abs{\theta}$
% for all $\theta\in \RR$, and letting
% \[
%     M = \int_E \abs{x}^m d\mu(x) < \infty    
% \]
% I claim that $\delta = \frac{\eps}{M}$ works for continuity. Indeed,
% for any $t'\in (t-\delta, t+\delta)$,
% \begin{align*}
%     \abs{f_m(t,x') - f_m(t,x)} &\leq \abs{\int \left(e^{it'x} - e^{itx}\right) x^m d\mu(x)} \\
%     &\leq \int \abs{e^{it'x} - e^{itx}} \abs{x}^m d\mu(x) \\
%     &\leq \int \abs{t' - t} \abs{x}^m d\mu(x) \\
%     &\leq \abs{t' - t} M \\
%     &< \eps
% \end{align*}
% (alternatively, fourier transforms are automatically continuous...)

% Problem 3.
% (i) Show that for any real numbers a, b one has ∫ b
% a eitxdx → 0 as |t| → ∞.
% (ii) Show that, for any f ∈ L1(R), the Fourier transform
% ̂
% f (t) =
% ∫ ∞
% −∞
% eitxf (x)dx
% tends to 0 as |t| → ∞.
\section{}
\begin{prob}
    Show that for any real numbers $a, b$ one has
    $\int_a^b e^{itx} dx \to 0$ as $\abs{t} \to \infty$.
\end{prob}

We can just compute:
\begin{align*}
    \abs{\int_a^b e^{itx} dx} &= \abs{\frac{e^{itx}}{it} \bigg|_a^b} \\
    &= \abs{\frac{e^{itb} - e^{ita}}{it}} \\
    &= \abs{\frac{e^{itb} - e^{ita}}{t}} \\
    &\leq \frac{2}{\abs{t}} \to 0
\end{align*}
where the $2$ is by triangle inequality or whatever.

\begin{prob}
    Show that, for any $f \in L^1(\RR)$, the Fourier
    transform $\hat{f}(t) = \int_{-\infty}^\infty e^{itx} f(x) dx$
    tends to $0$ as $\abs{t} \to \infty$.
\end{prob}
By splitting $f$ into its positive and negative components, WLOG $f$
is nonnegative and integral $\int_\RR f = M < \infty$.
By a well-known lemma, there exists an increasing
sequence of nonnegative, simple functions $f_n$ such that $f_n \to f$ pointwise;
fix this sequence; by monotone convergence theorem
or whatever you want,
\[
    \int_\RR f(x) - f_n(x) dx\to 0    
\]
In particular, for any $\eps$, I can find an $N$
such that $\int_\RR f(x) - f_N(x) < 0.0001\eps$.
% Now, for any $\eps$, I will find a $T$ such that if $\abs{t} > T$,
% then $\abs{\hat{f}(t)} < \eps$.

Then, for \emph{any} $t\in \RR$,
% Note that $\abs{e^{itx}(f(x) - f_n(x))} = f(x) - f_n(x)\leq f(x)$, thus
\[
    \abs{\int_\RR e^{itx}f(x) dx- \int_\RR e^{itx}f_N(x)dx} =\abs{\int_\RR e^{itx}(f(x) - f_N(x))dx} \leq \int_\RR f(x) - f_N(x) dx\leq 0.0001\eps
\]
% Well, pick an $N$ such that if $n \geq N$, the RHS is at most $0.0001\eps$.
Now, look carefully at
\[
    f_N = \sum_{k=1}^{K} c_k \id(E_k), \qquad \int_\RR f_N \leq \int_\RR f \implies \sum_k c_k\mu(E_k) \leq M
\]
In particular, each $E_k$ has finite, positive measure.
By Littlewood's first principle, for each $E_k$, there is a finite
collection of disjoint open intervals
$U_k = \bigcup_{n = 1}^{N_k}(a_{k,n}, b_{k,n})_{n=1}^{N_k}$ such that
\[
    \mu(E_k \Delta U_k) < 0.0001\eps / (K\cdot c_k)
\]
I.e. we can approximate $f_N$ with
\[
    g_N = \sum_{k=1}^{K} c_k \sum_{n=1}^{N_k} \id(a_{k,n}, b_{k,n})    
\]
Great, this means we get the following ridiculous bound:
\begin{align*}
    \abs{\int_\RR e^{itx}f_N - \int_\RR e^{itx}g_N}
    &\leq \abs{\sum_{k = 1}^K\int_\RR e^{itx} c_k \id(E_k) - \sum_{k=1}^K \int_\RR e^{itx} c_k \id(U_k)} \\
    &\leq \sum_{k = 1}^K\abs{\int_{E_k\setminus U_k} e^{itx}c_k - \int_{U_k\setminus E_k} e^{itx}c_k} \\
    &\leq \sum_{k = 1}^K\abs{\int_{E_k\setminus U_k} e^{itx}c_k} + \abs{\int_{U_k\setminus E_k} e^{itx}c_k} \\
    &\leq \sum_{k = 1}^K\abs{\int_{E_k\Delta U_k} e^{itx}c_k} \\
    &\leq \sum_{k = 1}^K \abs{c_k}\left(\frac{0.0001\eps}{K\cdot c_k}\right) \\
    &= 0.0001\eps
\end{align*}
Finally, note that
\[
    \int_\RR e^{itx}g_N = \sum_{k=1}^{K} c_k \sum_{n=1}^{N_k} \int_{a_{k,n}}^{b_{k,n}} e^{itx}    
\]
but every term of this summation can be made arbitrarily small by part (i).
In particular, this whole thing can be made smaller than $0.0001\eps$ in absolute value.
Observe $0.0001\eps + 0.0001\eps + 0.0001\eps \leq \eps$. Win.

% Problem 4. We say that f ∈ L2(R) is L2-differentiable with L2-derivative Df if
% ||τhf − f − hDf ||2/h → 0 as h → 0,
% where τhf (x) = f (x + h). Show that the function f (x) = max (1 − |x|, 0) is L2-differentiable and
% find its L2-derivative.
% Suppose that f ∈ L1 ∩ L2 is L2-differentiable. Show that u ̂f (u) ∈ L2. Deduce that f has a
% continuous version and that ||f ||∞ ≤ C||(1 + |u|) ̂f (u)||2 for some constant C < ∞, to be determined.
\section{}
\begin{prob}
    We say that $f \in L^2(\RR)$ is $L^2$-differentiable
    with $L^2$-derivative $Df$ if
    $\norm{\tau_h f - f - hDf}_2/h \to 0$ as $h \to 0$,
    where $\tau_h f(x) = f(x + h)$. Show that the
    function $f(x) = \max(1 - \abs{x}, 0)$ is
    $L^2$-differentiable and find its $L^2$-derivative.
\end{prob}

Well, we're basically allowed to ignore non-differentiable sets
of measure $0$. The $L^2$ derivative is
\[
    Df = \begin{cases}
        0 & x < -1\\
        1 & -1 \leq x < 0\\
        -1 & 0 \leq x < 1\\
        0 & x \geq 1
    \end{cases}    
\]
such that for all $\abs{h} < 1$,
\[
    \tau_hf - f - hDf = \begin{cases}
        0 & x < -1-h\\
        x+1+h & -1-h \leq x < -1 \\
        0 & -1 \leq x < -h\\
        -2x-2h & -h \leq x < 0\\
        0 & 0 \leq x < 1 - h\\
        x-1+h & 1 - h \leq x < 1\\
        0 & x \geq 1
    \end{cases}    
\]
Okay, the squared-$L^2$ norm of this thing (all parts
integrated with usual Riemann integral) is 
\[
    \frac{h^3}{3} + 4\cdot \frac{h^3}{3} + \frac{h^3}{3} = 2h^3
\]
Thus the $L^2$ norm is $\sqrt{2}h^{3/2}$, and the limit
is of $\sqrt{2}h^{1/2}$, which vanishes as $h \to 0$. Done.

\begin{prob}
    % [Extended]
    Suppose that $f \in L^1 \cap L^2$ is $L^2$-differentiable.
    Show that $u\hat{f}(u) \in L^2$.

    % In fact, everything works the way you want:
\end{prob}
Note that $Df \in L^2$, because
otherwise the above integral
is always $\infty$.

% Obviously, the main claim is that
% \[
%     \widehat{Df}(u) = iu\hat{f}(u)    
% \]
Well... let's just do it! Rewrite:
\[
    \frac{\tau_h f - f}{h}\stackrel{\norm{\bullet}_2}{\to} Df
\]
% Then, for all $x$,
% \[
%     \frac{\tau_h f(x) - f(x)}{h} e^{-iux} \stackrel{\norm{\bullet}_2}{\to}
%     Df(x) e^{-iux}
% \]
(as functions in $u$). We know that the Fourier transform
(or rather, its extension from $\LL^1\cap \LL^2$ to $\LL^2$)
is a Hilbert space isomorphism (theorem $7.6.1$),
so it preserves limits in the $L^2$ norm. Thus,
\[
    \int_\RR \frac{\tau_h f(x) - f(x)}{h} e^{iux} dx \stackrel{\norm{\bullet}_2}{\to} [\widehat{Df}]
\]
as a function in $u$.

\idea{
    [Sanity Check]
    As the main point of this problem is that $[\widehat{Df}]\in \LL^2$,
    let me talk about it.

    Here, we are defining $[\widehat{Df}]$ as described in $7.6.1$;
    it is the limit of the fourier transforms of a sequence
    of functions in $\LL^1\cap \LL^2$ which converge in $\LL^2$.
    As the Plancherel identity holds for all functions in
    $\LL^1\cap \LL^2$, we know that this limiting function is
    also in $\LL^2$.
}

However, observe that we can explicitly compute the LHS:
\begin{align*}
    \int_\RR \frac{\tau_h f(x) - f(x)}{h} e^{iux} dx
    &= \int_\RR f(x) e^{iux} \frac{e^{-ihu} - 1}{h} dx\\
    &= \frac{e^{-ihu} - 1}{h} \hat{f}(u)\\
    &\to -iu\hat{f}(u)
\end{align*}
\emph{pointwise}; this works \emph{literally everywhere}.
Convergence in $L^2$
implies that there exists a subsequence which converges
pointwise almost everywhere. Thus, $[\widehat{Df}] = [-iu\hat{f}]$,
and $u\hat{f}\in L^2$.

% Because $f\in L^1\cap L^2$, the Plancherel theorem applies,
% and we can write
% \[
%     \norm{\hat{f}}_2 = (2\pi)^{d/2}\norm{f}_2
% \]
% Thus $\hat{f}$ is automatically in $L^2$.

\begin{prob}
    Deduce that $f$ has a
    continuous version and that $\norm{f}_\infty \leq C \norm{(1 + \abs{u})\hat{f}(u)}_2$
    for some constant $C < \infty$, to be determined.
\end{prob}
Well,
\[
    [\hat{f}] = \left[\frac{i}{u} \widehat{Df}\right]    
\]
We know that $\widehat{Df}\in L^2$, and $\hat{f} \leq \norm{f}_1 < \infty$, thus
\begin{align*}
    \int_\RR \abs{\hat{f}(u)}du &= \int_{\abs{u}\leq 1} \abs{\hat{f}(u)} du + \int_{\abs{u}\geq 1}  \abs{\frac{\widehat{Df}(u)}{u}} du\\
    &\leq \sqrt{\int_{\abs{u}\leq 1} \abs{\hat{f}(u)}^2}\sqrt{\int_{\abs{u}\leq 1} 1} + \sqrt{\int_{\abs{u} \geq 1} \frac{1}{u^2} du}\cdot \sqrt{\int_{\abs{u}\geq 1}\abs{\widehat{Df}(u)}^2 du}\\
    &\leq \sqrt{2} \norm{\hat{f}}_2 + \sqrt{2}\norm{\widehat{Df}}_2\\
    &\leq \sqrt{2} \norm{\max{\hat{f}, \widehat{Df}}}_2 + \sqrt{2} \norm{\max{\hat{f}, \widehat{Df}}}_2\\
    &\leq 2\sqrt{2}\norm{(1 + \abs{u})\hat{f}(u)}_2
\end{align*}
Hence $\hat{f}\in L^1$, and the Fourier inversion theorem applies;
$f$ is a version of the inverse fourier transform of $\hat{f}$.

Thus, $f$ has a continuous version, because the inverse fourier
transform is continuous (for the same reasons as the fourier transform)
and $\norm{f}_\infty \leq \norm{\hat{f}}_1 / (2\pi)$ (for the same reasons
as the fourier transform). So you can pick $C = 100000$ or
whatever and we're done.

% \idea{
%     Note that if you split the integral on $\abs{u} < \norm{f}_1^{-1/2}$
%     instead, you get a $\norm{f}_1^{1/2} + \norm{\widehat{Df}}_2$ bound.
% }

% $f$ is also essentially-pointwise bounded by $\norm{\hat{f}}_1$.

% Now, the claim is that
% \[
%     \norm{f}_\infty \leq C_1\norm{f}_2 + C_2\norm{Df}_2  
% \]
% for some $C_1$ and $C_2$. WLOG $f$ is continuous.
% Scale $f$ such that $\norm{f}_\infty = 1$. Suppose now
% that $\norm{Df}_2 \leq 1$. We will show that $\norm{f}_2$ has
% a positive lower bound.


% Problem 5. Let (Xn : n ∈ N) be a sequence of random variables in R and let X be another random
% variable. Show that Xn → X weakly if and only if Xn → X in distribution.
\section{Portmanteau Theorem}
\begin{prob}
    Let $(X_n : n \in \NN)$ be a sequence of random variables
    in $\RR$ and let $X$ be another random variable.
    Show that $X_n \to X$ weakly if and only if
    $X_n \to X$ in distribution.
\end{prob}

This theorem is famous, important, and also completely tedious.

Nobody cares about your random variables, let $\mu_n$ and $\mu$
be their laws and we'll talk about those.

\subsection{Upgrading to Metric spaces}

Aiya first let me upgrade ``convergence in distribution" to
a more general form valid on metric spaces called ``convergence in pineapples"
so I can work with morally correct objects. The overarching claim for this section:

\begin{claim*}
    [Convergence in distribution $\iff$ Convergence in pineapples]
    Recall that a sequence of probability measure $\mu_n$ on $\RR$
    endowed with cumulative distribution function $F_n$
    \vocab{converges in distribution} to $\mu$ if
    \[
        F_n(x)\to F(x)   
    \]
    for all $x$ such that $F$ is continuous.

    On an arbitrary metric space $X$, a sequence of probability measures $\mu_n$
    on $X$ \vocab{converges in pineapples} to $\mu$ if
    for all open sets $U\subset X$,
    \[
        \lim\inf\mu_n(U) \geq \mu(U)
    \]

    Then, convergence in distribution $\iff$ convergence in pineapples.
\end{claim*}

First, let me delete the stupid CDFs from the definition.
We know $F(x) = \mu((-\infty, x])$,
and this is clearly already continuous from the right
(by continuity of measure from above in probability spaces).
\begin{claim*}
    I claim that $F(x)$ is continuous at $x$ iff $\mu(\{x\}) = 0$.
\end{claim*}
\begin{proof}
    The lower limit is
    \[
        \lim_{\eps\to 0} F(x - \eps) = \lim_{\eps\to 0} \mu((-\infty, x - \eps]) = \mu((-\infty, x)) = F(x) - \mu(\{x\})
    \]
    by continuity of measure from below. Thus the lower limit
    coincides with the upper limit of $F(x)$ iff $\mu(\{x\}) = 0$.
\end{proof}

Great. Let's show distribution $\implies$ pineapples now.

Recall that any open set $U\subset \RR$ is a \emph{countable}
disjoint union of open intervals $(a,b)$, and thus the problem reduces
to the case of a single open interval $(a,b)$.

But this is basically clear... here are four cases.
\begin{itemize}
\end{itemize}


% Thus, for instance,
% at least a cocountable number of points are continuous
% for all $X_n$ and $X$.

\begin{claim*}
    [Limit Interchange]
    Suppose you have a double sequence of real numbers $0\leq a_{mn}\leq 1$.
    Suppose also that $a_{mn}$ is non-decreasing in $m$ and non-decreasing in $n$;
    for \emph{all} $m \leq m'$ and $n \leq n'$, we have $a_{mn} \leq a_{m'n'}$.
    Then, we can swap the limits:
    \[
        \lim_{n\to \infty}\lim_{m\to \infty} a_{mn} = \lim_{m\to \infty}\lim_{n\to \infty} a_{mn}
    \]
\end{claim*}
\begin{proof}
    \begin{align*}
        \lim_{n\to \infty}\lim_{m\to\infty} a_{mn}
        &= \lim_{n\to \infty} \sup_{m} a_{mn} \\
        &= \sup_{n} \sup_{m} a_{mn} \\
        &= \sup_{m} \sup_{n} a_{mn} \\
        &= \sup_{m} \lim_{n\to \infty} a_{mn} \\
        &= \lim_{m\to \infty} \lim_{n\to \infty} a_{mn}
    \end{align*}
    We were able to set $\lim_n\sup_m = \sup_n\sup_m$
    because supremums preserve inequalities, and
    thus $\sup_{m} a_{mn}$ is increasing in $n$.
\end{proof}

Now, suppose $\mu_n\to \mu$ weakly. Pick some $a$ such that $F(x)$
is continuous. The first order of business
is to model two functions: $\id_{a < x}$, and $\id_{a \neq x}$.
Let:
\begin{itemize}
    \item $f_n$ be $1$ on $x\leq a - \frac1n$, $0$ on $x\geq a$, and linear in between.
    \item $g_n = \min(n\abs{x - a}, 1)$; this forms a $V$-shape near $a$.
\end{itemize}
Observe that $f_n$ and $g_n$ are continuous functions bounded between $0$ and $1$
which are increasing and converges pointwise to $\id_{x \leq a}$
and $\id_{x\neq a}$, respectively.
\idea{
    It might seem odd that we chose to model these instead of $\id_{a\leq x}$
    and $\id_{a = x}$, but I'm pretty sure we need to pass through the limit
    interchange lemma above, and I can only do this when things are increasing.
}

Seeing as this is a rather technical problem, let me justify each step.
% I guess I should justify these steps.
% \begin{enumerate}
%     \item Definition of $F$.
%     \item Definition of $f_n$.
%     \item Monotone convergence theorem.
%     \item Definition of weak convergence.
%     \item Claim.
%     \item Monotone convergence theorem.
%     \item Definition of $f_n$.
%     \item Definition of $F_m$.
% \end{enumerate}
\begin{align*}
    F(a) - \mu(\{a\}) &= \int \id_{x < a} d\mu \tag{by definition of $F$} \\
    &= \int \lim_{n\to \infty} f_n d\mu \tag{by definition of $f_n$} \\
    &= \lim_{n\to \infty} \int f_n d\mu \tag{by monotone convergence theorem} \\
    &= \lim_{n\to \infty} \lim_{m\to \infty} \int f_n d\mu_m \tag{by definition of weak convergence} \\
    &= \lim_{m\to \infty} \lim_{n\to \infty} \int f_n d\mu_m \tag{by limit interchange} \\
    &= \lim_{m\to \infty} \int \lim_{n\to \infty} f_n d\mu_m \tag{by monotone convergence theorem} \\
    &= \lim_{m\to \infty} \int \id_{x < a} d\mu_m \tag{by definition of $f_n$} \\
    &= \lim_{m\to \infty} F_m(a) - \mu_m(\{a\}) \tag{by definition of $F_m$} \\
\end{align*}

Running this exact same calculation on $g_n$ shows
\[
    1 - \mu_m(\{a\}) = \lim_{m\to \infty} 1 - \mu_m(\{a\})    
\]
Subtracting these two results, we find that $F_m(a)\to F(a)$
for all $a$ such that $F$ is continuous. This concludes.

\idea{
    Sheesh, thread the needle.
}

Now, suppose $X_n \to X$ in distribution. Give me a continuous bounded
function $f$. I'm going to approximate $f$ by a series of piecewise
constant functions, using standard real analysis.

\begin{psec*}
    [Constructing $f_n$]
    Here's how we construct function $f_n$. First, find some $M_n$
    such that $\mu(\abs{X} \geq M_n) < \frac{1}{n\norm{f}_\infty}$ (this is possible because
    $\mu$ is a probability measure).

    Set $f_n = 0$ outside $[-M_n, M_n]$. Second, on $[-M_n, M_n]$,
    $f$ is continuous, and
    thus uniformly continuous (by standard results). Let $\eps = \frac{1}{n}$.
    %  and thus
    % can be approximated pointwise within $\eps = \frac1n$
    % by a piecewise constant function $f_n$. This is also a standard
    % result, but I want to explicitly show it is bounded, so let me
    % just do it:
    \begin{itemize}
        \item $f$ is UC, thus $\forall\eps\exists \delta$
        such that if $\abs{x - y} < \delta$, then $\abs{f(x) - f(y)} < \eps$.
        \item Use $\delta$ as a discretization; split $[a,b]$
        into a finite number of intervals of length at most $\delta$, 
        taking care to avoid a countable number of points which will
        be specified shortly.
        \item In each interval $I_i$, set $f_k(I_i)$ to be an arbtrarily chosen
        value from $f(I_i)$.
        \item In particular, $f_k$ is uniformly bounded by
        the same bounds as $f$.
    \end{itemize}
    Then, $\abs{f_n} < \norm{f}_\infty$ everywhere, and $\abs{f(x) - f_n(x)} < \eps$
    everywhere on $[-M_n, M_n]$.
\end{psec*}


Finally, give me some $\eps$, and I'll find you an $n$ such that
$\abs{\mu(f) - \mu_m(f)} < \eps$ for all $m > n$.

First, let $k > \frac{0.0001}{\eps}$.
Find $n$ such that for all $m > n$,
\begin{itemize}
    \item $\mu_m((-\infty, M_k])$ is within $0.0001\eps/\norm{f}_\infty$ of $\mu((-\infty, M_k])$.
    \item $\mu_m((-\infty, -M_k])$ is within $0.0001\eps/\norm{f}_\infty$ of $\mu((-\infty, -M_k])$.
    \item $\mu_m(f_k) - \mu(f_k) < 0.0001\eps$.
\end{itemize}
This is possible because $\mu_m \to \mu$ weakly. (The last line is true
because $\mu_m(f_k)$ is just the measure of a bunch of intervals).

\idea{
    [Detail]
    Okay, we only know $\mu_m((-\infty, x])\to \mu((-\infty, x])$
    for all $x$ such that $F$ is continuous.

    However:
    \begin{itemize}
        \item $F$ is discontinuous at $x$ iff $\mu(\{x\}) > 0$.
        \item Suppose there existed an uncountable set $S$ such that
        $\mu(\{x\}) > 0$ for all $x\in S$. Then, one can choose a finite
        subset of $S$ with arbitrarily large measure (by splitting $S$ into
        buckets $(\frac1{n+1}, \frac1n)$ and finding a bucket with uncountably many things).
        This is a contradiction, thus the set of discontinuities is countable.
        \item But $\RR$ is uncountable.
    \end{itemize}
    So the upshot is that if $M_k$ fails, you can just poke it to a slightly
    larger value for which $F$ is actually continuous, and ignore this problem.

    And actually, the number of discontinous points for all $F_m$ is countable too.
    So in our definition of $f_k$ above, let's just pick our $I_i$'s to avoid
    all of these points.
}
%     and $\mu((-\infty, -M_k])$,
% for all $m > n$. 
% \end{itemize}

Then, for all $m > n$,
\begin{align*}
    \abs{\mu_m(f) - \mu_m(f_k)} &\leq \int \abs{f - f_k} \mu_m(dx)\\
    &\leq \int \abs{f}\id\left((-\infty, -M_k)\cup (M_k, \infty)\right) \mu_m(dx) + \int 0.0001\eps \id\left([-M_k, M_k]\right)\mu_m(dx)\\
    &\leq \norm{f}_\infty \mu_m((-\infty, -M_k] \cup [M_k, \infty))  + 0.0001\eps\\ % + 0.0001\eps\mu_m((-\infty, M_k]\cup (-M_k, \infty))
    &\leq \norm{f}_\infty \mu((-\infty, -M_k] \cup [M_k, \infty)) + \norm{f}_\infty \cdot \frac{0.0001\eps}{\norm{f}_\infty} + 0.0001\eps\\
    % &+ 0.0001\eps\mu_m((-\infty, M_k]\cup (-M_k, \infty)) + 0.0001\eps\\
    &\leq 0.0001\eps + 0.0001\eps + 0.0001\eps\\
    &\leq 0.0003\eps 
\end{align*}
% Similarly, $\abs{\mu(f) - \mu(f_k)}\leq 0.0003\eps$.
Similarly,
\begin{align*}
    \abs{\mu(f) - \mu(f_k)} &\leq \int \abs{f - f_k} \mu(dx)\\
    &\leq \int \abs{f}\id\left((-\infty, -M_k)\cup (M_k, \infty)\right) \mu(dx) + \int 0.0001\eps \id\left([-M_k, M_k]\right)\mu(dx)\\
    &\leq \norm{f}_\infty \mu((-\infty, -M_k] \cup [M_k, \infty))  + 0.0001\eps\\ % + 0.0001\eps\mu_m((-\infty, M_k]\cup (-M_k, \infty))
    &\leq 0.0001\eps + 0.0001\eps\\
    &\leq 0.0002\eps
\end{align*}
Finally, the difference between $\mu_m(f_k)$ and $\mu(f_k)$ is at most $0.0001\eps$ from before.
So in total $\abs{\mu_m(f) - \mu(f)} < 0.0006\eps$. Done.

% Thus, each $f_n$ is bounded and converges pointwise to $f$.
% Now, by dominated convergence theorem,
% \[
%     \int f_n d\mu \to \int f d\mu
% \]

% Problem 6. Let X1, . . . , Xn be independent N (0, 1) random variables. Show that
% (
% X,
% n∑
% m=1
% (Xm − X)2
% )
% and
% (
% Xn
% √n ,
% n−1∑
% m=1
% X2
% m
% )
% have the same distribution, where X = (X1 + · · · + Xn)/n.
\section{}
\begin{prob}
    Let $X_1, \dots, X_n$ be independent $N(0, 1)$ random variables.
    Show that
    % \[
    %     \left(
    %         \overline{X}, \frac{1}{n} \sum_{m=1}^n (X_m - \overline{X})^2
    %     \right)
    %     \text{ and }
    %     \left(
    %         \frac{X_n}{\sqrt{n}}, \frac{1}{n-1} \sum_{m=1}^n X_m^2
    %     \right)
    % \]
    \[
        \left(\overline{X}, \sum_{m = 1}^n (X_m - \overline{X})^2\right) \text{ and } \left(\frac{X_n}{\sqrt{n}}, \sum_{m=1}^{n-1} X_m^2\right) \tag{$\star$}
    \]
    have the same distribution, where $\overline{X} = (X_1 + \dots + X_n)/n$.
\end{prob}

First of all, here's a completely trivial fact:
\begin{claim*}
    If $X,Y:\Omega\to E$ are two random variables which have the same distribution
    and $f:E\to F$ is a measurable function, then $f(X)$ and $f(Y)$
    have the same distribution.
\end{claim*}
\begin{proof}
    The definition of having the same distribution is that
    \[
        \mu\circ X^{-1} = \mu\circ Y^{-1}
    \]
    Then clearly
    \[
        \mu\circ (f\circ X)^{-1} = \mu\circ X^{-1}\circ f^{-1} = \mu\circ Y^{-1}\circ f^{-1} = \mu\circ (f\circ Y)^{-1}
    \]
\end{proof}


Now, let me write down an orthogonal matrix $U$. Actually, I'm only
going to fill in the bottom row:
\[
    U = \begin{pmatrix}
            ? & \dots & ?\\
            \vdots & \ddots & \vdots\\
            ? & \dots & ?\\
            \frac{1}{\sqrt{n}} & \dots & \frac{1}{\sqrt{n}}
        \end{pmatrix}  
\]
and by e.g. Gram-Schmidt, I can fill in the rest of the matrix
with an orthonormal basis for the orthogonal
complement of $\vec{a} = \frac{1}{\sqrt{n}}(1,\dots, 1)^\top$.

% by telling you an orthonormal basis.
% \begin{itemize}
%     \item I want $U$ to send $\vec{a} = \frac{1}{\sqrt{n}}(1,\dots, 1)^\top$ to $(0,\dots,0,1)^\top$.
%     % such that
%     \item By e.g. Gram-Schmidt, I can find an orthonormal basis
%     for the orthogonal complement of $\vec{a}$. Pick one arbitrarily.
%     Send this basis to the remaining $n-1$ standard basis vectors.
% \end{itemize}
% This way,
% \[
%     \left(U\begin{pmatrix}
%         x_1\\
%         \vdots\\
%         x_n
%     \end{pmatrix}\right)_n = \frac{x_1 + \dots + x_n}{\sqrt{n}}
% \]
% \begin{itemize}
%     \item I want $U$ to send $(0,\dots,0,1)^\top$ to $\vec{a} = \frac{1}{\sqrt{n}}(1,\dots, 1)^\top$.
%     \item By e.g. Gram-Schmidt, I can find an orthonormal basis
%     for the orthogonal complement of $\vec{a}$. Pick one arbitrarily.
%     Send the remaining $n-1$ standard basis vectors to this basis.
% \end{itemize}

Great! Consider the function
\[
    f:\RR^n \to \RR^2: \begin{pmatrix}
        x_1\\
        \vdots\\
        x_n
    \end{pmatrix}
    \to \begin{pmatrix}
        \frac{x_n}{\sqrt{n}}\\
        \norm{x}_2^2 - x_n^2
    \end{pmatrix}    
\]

It is clear that the RHS of $\star$ is $f(X)$
where $X\in \RR^n = (X_1, \dots, X_n)^\top$.
I claim that the LHS is equal to $f(UX)$,
where $UX$ is a completely legal thing to write down
(I am composing $U\circ X$ where $U:\RR^n\to \RR^n$).

It suffices to show that (as linear algebra identities)
\[
    \sum_{m = 1}^n (X_m - \overline{X})^2 = \norm{UX}_2^2 - (UX)_n^2 = \norm{X}_2^2 - n\overline{X}^2
\]
This is basically obvious and left as an exercise. No I'm kidding,
uh, the LHS is
\[
    \sum X_i^2 - 2\overline{X}\sum X_i + \overline{X}^2 \cdot n = \norm{X}_2^2 - n\overline{X}^2
\]
lol.
% No, I'm kidding. Look, the LHS is just the norm-square of
% \[
%     Y = \begin{pmatrix}
%         X_1 - \overline{X}\\
%         \vdots\\
%         X_n - \overline{X}
%     \end{pmatrix}
%     =
%     \underbrace{\begin{pmatrix}
%         \frac{n-1}{n} & -\frac1n & \dots & -\frac1n\\
%         -\frac1n & \frac{n-1}{n} & \dots & -\frac1n\\
%         \vdots & \vdots & \ddots & \vdots\\
%         -\frac1n & -\frac1n & \dots & \frac{n-1}{n}
%     \end{pmatrix}}_Q
%     X
% \]
% But it's clear that $I_n-Q = P$ is the orthogonal
% projector onto $\vec{a}$, and thus $Q$ is the orthogonal
% projector onto the orthogonal complement of $\vec{a}$.
% And also, $n\overline{X}^2 = \norm{PX}_2^2$. So we're just saying
% \[
%     \norm{QX}_2^2 = \norm{X}_2^2 + \norm{PX}_2^2
% \]
% but this is just the Pythagorean theorem.

% Okay win!
By the claim, it suffices to show that
$UX$ and $X$ have the same distribution.
But if you stare at their distributions, this is completely obvious:
\[
    f_{UX}(\vec{x}) = \frac{1}{(2\pi)^{n/2}} e^{-\frac{1}{2}\norm{U\vec{x}}_2^2} = \frac{1}{(2\pi)^{n/2}} e^{-\frac{1}{2}\norm{\vec{x}}_2^2} = f_X(\vec{x})
\]

% Note that $U$ is orthogonal so $\norm{UX}_2^2 = \norm{X}_2^2$,
% and $(UX)_n^2 = \frac{1}{n}\overline{X}$.


% By the claim, it suffices to prove that
% \[
%     \begin{pmatrix}
%         X_1 - \overline{X}\\
%         \vdots\\
%         X_n - \overline{X}\\
%         \overline{X}\sqrt{n}
%     \end{pmatrix}
%     \text{ and }
%     \begin{pmatrix}
%         X_1\\
%         \vdots\\
%         X_n
%     \end{pmatrix}    
% \]
% have the same distribution.

% Well, I'll prove that
% \[
%     \left(\overline{X}, (X_m - \overline{X})^2\right) \text{ and } \left(\frac{X_n}{\sqrt{n}}, \sum_{m=1}^{n-1} X_m^2\right) 
% \]


% Well, $X = (X_1,\dots, X_n)^\top\in \RR^n$ is a jointy Gaussian random variable.
% Let $u = \frac1n(1,1,1\dots)^\top$.
% We can write
% \[
%     \EE e^{i \overline{X}} = \EE e^{i \braket{u | X}} = \EE \prod_k e^{iu_kX_k} = e^{-\abs{u}^2/2}
% \]
% Thus $\overline{X}$ is $N(0, 1/n)$. Of course, $\frac{X_n}{\sqrt{n}}\sim N(0, 1/n)$
% as well.

% \begin{prob}
%     Show that
%     \[
%         \sum_{m=1}^n (X_m - \overline{X})^2
%         \text{ and }
%         \sum_{m=1}^{n-1} X_m^2
%     \]
%     have the same distribution.
% \end{prob}

% The RHS is the norm-square of an RV with distribution $N(0,I_{n-1})$.
% The LHS is the norm-square of
% \[
%     Y = \begin{pmatrix}
%         X_1 - \overline{X}\\
%         \vdots\\
%         X_n - \overline{X}
%     \end{pmatrix}
%     =
%     \underbrace{\begin{pmatrix}
%         \frac{n-1}{n} & -\frac1n & \dots & -\frac1n\\
%         -\frac1n & \frac{n-1}{n} & \dots & -\frac1n\\
%         \vdots & \vdots & \ddots & \vdots\\
%         -\frac1n & -\frac1n & \dots & \frac{n-1}{n}
%     \end{pmatrix}}_V
%     X
% \]
% which is jointly Gaussian (by theorem $8.2.1$ or common sense).
% Clearly it suffices to show that (in a suitable sense)
% $Y\sim N(0, I_{n-1})$ as well.

% Please observe that $I - V$ is merely the (orthogonal) projection
% onto the subspace spanned by $(1,\dots, 1)^\top$; thus $V$
% is the orthogonal projection onto the $n-1$-dimensional subspace
% orthogonal to $(1,\dots, 1)^\top$. In particular, $V$ is
% orthogonally equivalent to any other orthogonal projection on
% to an $n-1$-dimensional subspace, such as
% \[
%     P_{n-1} = \begin{pmatrix}
%         1 & 0 & \dots & 0 & 0\\
%         0 & 1 & \dots & 0 & 0\\
%         \vdots & \vdots & \ddots & \vdots & \vdots\\
%         0 & 0 & \dots & 1 & 0\\
%         0 & 0 & \dots & 0 & 0
%     \end{pmatrix}
% \]

% From the proof of part $(d)$ in theorem $8.2.1$ (and/or common sense)
% I know that $\var Y = \EE[V XX^\top V^\top ] = VV^\top$ and $\EE[Y] = 0$.
% But $V$ is just some projection, thus $VV^\top$ is orthogonally
% equivalent to some $n-1$ dimensional projection as above.
% The distributions of Gaussians are fixed by $\var$ and $\EE$,
% thus $Y$ has the same distribution as
% \[
%     P_{n-1}X = \begin{pmatrix}
%         X_1\\
%         \vdots\\
%         X_{n-1}\\
%         0
%     \end{pmatrix}
% \]
% To be exact, let $Q_{n-1} : \RR^n\to \RR^{n-1}$ be
% an actual projection matrix.
% Thus note that $\norm{X} = \norm{Q_{n-1}X}$,
% and $Q_{n-1}X \sim N(0, I_{n-1})$ as desired.

% Thus, the question is to show that if $X\sim N(0,I_n)$,
% and $Y\sim N(0, I_{n-1})$, then $n\norm{X}^2$ is distributed
% the same as $(n-1)\norm{Y}^2$.

% Okay I guess we just have to compute at this point.
% The distribution of $\norm{X}^2$ is


% Then, $Y$ is also joinly Gaussian, with
% \[
%     \mu = \EE[Y] = \EE[X] - \left(\frac{\EE[X_1] + \dots + \EE[X_n]}{n}\right)\begin{pmatrix}
%     1 \\ \vdots \\ 1\end{pmatrix} = 0
% \]
% Observe that the RHS is precisely $\frac1n$ of
% the norm of this random variable.


% We know that the distribution of a sum is simply the convolution
% of the distributions.


% Thus, the distribution of $\overline{X}$
% is
% \begin{align*}



% Problem 7. Let (E, E, μ) be a measure space and τ : E → E a measure-preserving transfor-
% mation. Show that Eτ = {A ∈ E : τ −1(A) = A} is a σ-algebra, and that a measurable function f is
% Eτ -measurable if and only if it is invariant, that is f ◦ τ = f .
\section{}
\begin{prob}
    Let $(E, \mathcal{E}, \mu)$ be a measure space and
    $\theta : E \to E$ a measure-preserving transformation.
    Show that $\mathcal{E}_\theta = \{A \in \mathcal{E} : \theta^{-1}(A) = A\}$
    is a $\sigma$-algebra.
\end{prob}

$E\subset \mathcal{E}_\theta$ because $\theta^{-1}(E) = E$. Suppose $A\in \mathcal{E}_\theta$; then
\[\theta^{-1}(E\setminus A) = \theta^{-1}(E)\setminus \theta^{-1}(A) = E\setminus A\implies E\setminus A\in \mathcal{E}_\theta.\]
Finally, if $A_1,\dots A_n\in \mathcal{E}_\theta$, then
\[
	\theta^{-1}\left(\bigcup_{i} A_i\right) = \bigcup_{i}\theta^{-1}\left(A_i\right) = \bigcup_{i} A_i
\]
as desired.

\begin{prob}
    Show that a measurable function
    $f$ is $\mathcal{E}_\theta$-measurable if and only if it is invariant,
    that is $f \circ \theta = f$.
\end{prob}
I will assume $f$ points into $\RR$. Otherwise, there are counterexamples;
the measurability condition can't distinguish between points if the measure
space isn't ``Hausdorff''.

% Suppose $f$ maps into a space $(F, \mathcal{F})$.
Suppose $f$ is
$\mathcal{E}_\theta$-measurable. Then, noting that points are measurable
in $\RR$,
$f^{-1}(x)$ is measurable in $\mathcal{E}_\theta$. This implies
$\theta^{-1}(f^{-1}(x)) = f^{-1}(x)$ for all $x$. Now, give
me any $x\in E$. I know that $x\in \theta^{-1}(f^{-1}(f(x)))$.
Thus,
\[
    f(\theta(x)) = f(f^{-1}(f(x))) = f(x)    
\]

On the other hand, suppose $f\circ\theta = f$. Then, for all measurable
sets $A\subset \RR$, $f^{-1}(A) = \theta^{-1}(f^{-1}(A))$,
i.e. $f^{-1}(A) \in \mathcal{E}_\theta$.



% Problem 8. Show that, if θ is an ergodic measure-preserving transformation and f is a θ-invariant
% function, then there exists a constant c ∈ R such that f = c a.e.
\section{}
\begin{prob}
    Show that, if $\theta$ is an ergodic measure-preserving
    transformation and $f$ is a $\theta$-invariant
    function, then there exists a constant $c \in \RR$ such
    that $f = c$ a.e.
\end{prob}
If $E$ has measure $0$, the problem is stupid, because all functions
are equal almost everywhere. Suppose otherwise.

Observe that for any measurable $A\subset \RR$, $f^{-1}(A)$ is $\theta$-invariant;
by ergodicity, either $f^{-1}(A)$ or $E\setminus f^{-1}(A)$ has measure $0$.

Now, consider the sets $E_i = f^{-1}([i, i+1))$ for $i\in \ZZ$;
this is a countable set, thus they cannot \emph{all} have measure $0$.
Two of them cannot have nonzero measure, else neither of their complements
would have measure $0$. Thus, exactly
one of them has nonzero measure; call it $k = i$.

Then, consider the dyadic blocks
\[
    E_{n,a} = \left[\frac{a}{2^n} + k, \frac{a+1}{2^n} + k\right), \qquad a = 0, \dots, 2^n - 1
\]
These have the property that for all $n$, they partition $[k,k+1)$.
Thus, by the above reasoning, for each $n$ there is exactly one $a$
such that $E_{n,a}$ has nonzero measure. Call this $a_n$.
Clearly $a_{n+1} \in \{2a_n, 2a_n + 1\}$; any other choice would contradict
the fact that $E\setminus E_{n,a}$ has zero measure.

Thus, we find ourselves with a descending series of blocks
\[
    E_{0,a_0} \supset E_{1,a_1} \supset \dots    
\]
such that $E\setminus E_{k, a_k}$ has measure $0$. By
continuity of measure from below,
\[
    \mu\left(E\setminus \inf E_{k, a_k}\right) = \lim_{k\to \infty} \mu\left(E\setminus E_{k, a_k}\right) = 0   
\]
But clearly $\inf E_{k,a_k}$ is some singleton $\{c\}$ (by definition of the real numbers).
Thus, $E\setminus f^{-1}(\{c\})$ has measure $0$, and we are done.



\end{document}